\providecommand{\PushK}{890666}
\providecommand{\PushTheta}{0.99999999997327}

\section{The whole argument with one query}
\label{sec:mechanism}

\paragraph{What this section proves and how to read it.}
This section proves nothing new; it assembles, with one query instead
of six, the four facts on which the whole proof rests, in the form
\cref{prop:mech:one-query}: at a least counterexample,
$1\le\mathcal E+BR_I/2+W_1+o(1)$, where
$W_1=BR_I/(2\Kint^{35/36})$ is the one-query wrapping term. The
four facts are: (1) a
reflection-positive law gives $\alpha^{72}\le Y$ for the $72$-fold
Fourier average $Y$ (\cref{subsec:mech:chessboard}); (2) at a least
counterexample every change of the law is bounded by its price times
$Y$, because fixing labels restricts to shorter square-step
progressions (\cref{subsec:mech:prices}); (3) small primes enter
through an external index coupling the two ends of the query
(\cref{subsec:mech:index}); (4) coupled terms in which the query is
accepted and intervals do not wrap are negligible by the kernel
estimate (\cref{subsec:mech:kernel}). Inputs: the definitions of
Part~II, quoted by label. Output: the one-query inequality, whose
six-query form is Theorem~C (\cref{thm:comparison}); Part~V applies
it with six queries. A first reading should follow the statements and
the sketch after \cref{prop:mech:one-query} and skip the quoted
hypothesis checks; nothing here is a premise of a later section.
\Cref{subsec:mech:budget} explains why one query cannot give
$\Theta<1$, and \cref{subsec:mech:crude} why conditioning independent
labels does not work, which motivates \cref{sec:design-walk}.

\subsection{The \texorpdfstring{$72$}{72}-fold average and reflection positivity}\label{subsec:mech:chessboard}

Let $N$ be large (in particular so large that the prime ranges
below are nonempty as stated), let $A\subseteq[N]$ be square-difference-free, and let
$\alpha=|A|/N$. The labels of the proof live in the finite group
$G_N=\prod_p\Z/p^{e_p}\Z$, a product over three kinds of primes
(\cref{def:record:functional}): the primes $p\equiv1\pmod4$ below
$\Pfour=2^{512}$, with $e_p=2\mcl$ and $\mcl=8192$; the primes
$p\equiv3\pmod4$ in $[Q_3,\Pfour)$, $Q_3=19951586$; and all primes in
$[\Pfour,N^{1/20}]$, the last two kinds with $e_p=1$. The prime $2$ and the
primes $p\equiv3\pmod4$ below $Q_3$ are not coordinates of $G_N$; they enter
in \cref{subsec:mech:index}. The \emph{lift} of $A$ is
\[
 F_A(z)=\sum_{q(\xi)\le N^{1/20}}\widehat{\one_A}(\xi)\,e(z\cdot\xi),
 \qquad\widehat{\one_A}(\xi)=\frac1N\sum_{n\in A}e(-n\xi),
\]
the sum over the characters $\xi$ of $G_N$ whose conductor $q(\xi)$ is at
most $N^{1/20}$. It is real, and its mean over the uniform law on $G_N$ is the
coefficient of the trivial
character: $\E F_A=\alpha$.

\paragraph{Slots, cuts and reflection positivity.}
Let $V$ be a finite set of \emph{slots}, $d=|V|$. For a finite signed
measure $\kappa$ on the labellings $z=(z_v)_{v\in V}\in G_N^V$ and real
functions $g_v$ on $G_N$ put
\[
 \Lambda_\kappa\bigl((g_v)_{v\in V}\bigr)=\int\prod_{v\in V}g_v(z_v)\,d\kappa(z),
 \qquad\Lambda_\kappa(g)=\Lambda_\kappa\bigl((g)_{v\in V}\bigr).
\]
The average of the proof is $Y=\Lambda_\mu(F_A)$ for one fixed probability
measure $\mu$, the \emph{reference}. A \emph{cut} is a fixed-point-free
involution $\vartheta$ of $V$ with a decomposition $V=V_+\sqcup V_-$,
$V_-=\vartheta V_+$. The \emph{cut matrix} of $\kappa$ at $\vartheta$ is
$C_\vartheta(\kappa)_{x,y}=\kappa(z_{V_+}=x,\ z_{\vartheta V_+}=y)$, indexed by
$x,y\in G_N^{V_+}$, and its \emph{cut form} is
$Q_\kappa(X,X')=\ip X{C_\vartheta(\kappa)X'}$ for \emph{half-functions}
$X,X'\colon G_N^{V_+}\to\R$. If $X$ and $X'$ are products of functions of the
single labels, $Q_\kappa(X,X')$ is the average $\Lambda_\kappa$ in which the
factors of $X$ sit on $V_+$ and those of $X'$ are carried to $V_-$ by
$\vartheta$ \eqref{eq:cc:bilinear}. The measure $\kappa$ is \emph{reflection
positive} at $\vartheta$ if $C_\vartheta(\kappa)$ is symmetric and positive
semidefinite (\cref{def:rp,def:cc:order}). Reflection positivity at a cut is
preserved by sums, positive multiples and tensor products
\eqref{eq:cc:kron}.

\paragraph{One cut, and folding repeatedly until all slots agree.}
Let $\phi_+$ be the map of $V$ that fixes $V_+$ and sends $v\in V_-$ to
$\vartheta v$, and $\phi_-$ the map that fixes $V_-$ and sends $v\in V_+$ to
$\vartheta v$. If $\kappa$ is reflection positive at $\vartheta$, the
Cauchy--Schwarz inequality for its cut form gives, for every assignment
$h\colon V\to\R^{G_N}$ of functions to slots (\cref{lem:holder:one-cut}),
\begin{equation}\label{eq:mech:one-cut}
 \bigl|\Lambda_\kappa(h)\bigr|^2\le\Lambda_\kappa(h\circ\phi_+)\,
 \Lambda_\kappa(h\circ\phi_-),\qquad\Lambda_\kappa(h\circ\phi_\pm)\ge0,
\end{equation}
where $\Lambda_\kappa(h)=\Lambda_\kappa((h(v))_v)$: a mixed average is at most the
geometric mean of the two averages in which one half has been copied onto
the other. A sequence of cuts $\vartheta^1,\dots,\vartheta^T$ with signs
$s_t\in\{\pm\}$ is \emph{sorting} if the composition
$\phi^{\vartheta^T}_{s_T}\circ\dots\circ\phi^{\vartheta^1}_{s_1}$ of its folds is
constant (\cref{def:sorting}).

\begin{theorem*}[Signed H\"older inequality]
If $\kappa$ is reflection positive at every cut of a sorting sequence, then
for all real functions $g$, $g_v$ on $G_N$ \textup{(\cref{thm:holder})}
\[
 \Lambda_\kappa(g)\ge0\qquad\text{and}\qquad
 \bigl|\Lambda_\kappa\bigl((g_v)_{v\in V}\bigr)\bigr|
 \le\prod_{v\in V}\Lambda_\kappa(g_v)^{1/d}.
\]
\end{theorem*}

\emph{Sketch.} Fix finitely many functions $f_1,\dots,f_k$, and let $M$ be the
largest value of $|\Lambda_\kappa(h)|$ over the assignments $h$ with values in
$\{f_i\}$. If $h$ attains $M$, both folded assignments $h\circ\phi_\pm$ again
take values in $\{f_i\}$. By \eqref{eq:mech:one-cut} they are nonnegative,
at most $M$, and their product is at least $M^2$, so both equal $M$: folds of
maximizers are maximizers. Folding along the whole sorting sequence
produces a constant maximizer, so $M=\max_i\Lambda_\kappa(f_i)$. Applied to
$f_v=g_v/(\Lambda_\kappa(g_v)+\epsilon)^{1/d}$, $\epsilon>0$, multilinearity gives
the inequality as $\epsilon\to0$. This is the chessboard estimate of statistical mechanics in
finite form. Consequently $\|g\|_\kappa=\Lambda_\kappa(g)^{1/d}$ is a seminorm, and
if every one-slot marginal of $\kappa$ is uniform, then $|\E g|\le\|g\|_\kappa$,
by the inequality with $g$ at one slot and $1$ at the others
(\cref{cor:seminorm}). With $g=F_A$ this is $\alpha^d\le\Lambda_\kappa(F_A)$. The
inequality says more: every average of the same shape, with arbitrary real
functions at the slots, is controlled by the diagonal averages of its
inputs. This is what makes the prices of \cref{subsec:mech:prices} possible.

\paragraph{The slot array.}
The slots form $V=\Zsix\times\Zsix\times\{0,1\}$, $d=72$, with slots
$(u_0,u_1,\ell)$ (\cref{def:slot-array}). On each chamber coordinate $\Zsix$
there are the chamber reflections $\vartheta_a(j)=2a-1-j$ with plus halves
$H_a=\{a,a+1,a+2\}$, and $\vartheta_{c,a}$ is $\vartheta_a$ acting on $u_c$. The
layer carries the reserved cut $\iota(u_0,u_1,\ell)=(u_0,u_1,1-\ell)$, with plus
half $V_+=\{\ell=0\}$. The plus folds of $\vartheta_0$ and $\vartheta_4$ are
$f_0=(0,1,2,2,1,0)$ and $f_4=(0,0,5,4,4,5)$ (values at $0,\dots,5$), and
$f_0\circ f_4\circ f_0$ is constant. Folds of different coordinates commute,
so the seven-step sequence
$\vartheta_{0,0},\vartheta_{0,4},\vartheta_{0,0},\vartheta_{1,0},\vartheta_{1,4},
\vartheta_{1,0},\iota$, all with sign $+$, maps every slot to $(0,0,0)$
(\cref{lem:slot-array}(ii),(iii)). The family $\mathfrak F$ of the chamber
cuts and $\iota$ consists of seven involutions. The reference is a product
$\mu=\bigotimes_\Blk\mu_\Blk$ over the \emph{atoms} $\Blk$ of $G_N$ (the single
primes of the prefix, blocks of the primes $p\equiv3\pmod4$ in
$[Q_3,\Pfour)$, and the single tail primes) of probability measures that are
reflection positive at every cut of $\mathfrak F$ and invariant under the maps
$z_{v,p}\mapsto r_pz_{v,p}+t_p$, $r_p$ a square unit, chosen independently at
each prime; in particular their one-slot marginals are uniform
(\cref{hyp:record:reference}). Hence (\cref{lem:record:holder})
\begin{equation}\label{eq:mech:density}
 \alpha^{72}\le\Yf(\Lsc,A)=\Lambda_\mu(F_A)=Y .
\end{equation}
Positivity turns density into size; the proof shows, conversely, that
$Y\le C\Lsc^{-a}\alpha$ at the scale $\Lsc$ in the decomposition
$N=\Kint\Lsc$ of \cref{subsec:mech:prices} (the next subsection).

\paragraph{The six distinguished slot pairs (queries).}
The proof uses the six queries $e=(u_e,v_e)=((0,j,0),(2,\rot j,1))$,
$j\in\Zsix$, of the family $\mathcal E_6$, where $\rot j=j+2$ for even $j$ and
$\rot j=j-2$ for odd $j$. Their twelve
endpoints are distinct, the sources lie in $V_+$, the targets in $V_-$, and
the reflected targets $\iota v_e=(2,\rot j,0)$ in $V_+$
(\cref{lem:slot-array}(iv)). Each query crosses the reserved cut. In this
section one query $e$ is enough.

\subsection{Rescaling by square steps: the price of changing a local law}\label{subsec:mech:prices}

\paragraph{The induction.}
Write $N=\Kint\Lsc$ with a fixed integer $\Kint\ge2$ (in the
baseline $\Kint=\PushK$; the improved stage uses
$\Kint=\FourRoutesK$), and fix a price exponent $a$ with
$0<a<\sigma=1/825$,
where $\sigma$ is the error exponent of the analytic stage
(\cref{def:record:analytic-stage}). The proof shows that for a suitable
$C_0\ge1$
\[
 \Yf(\Lsc,A)\le C_0\Lsc^{-a}\alpha
\]
for every $\Lsc$ and every square-difference-free $A\subseteq[\Kint\Lsc]$. With
\eqref{eq:mech:density} this gives $\alpha^{71}\le C_0\Lsc^{-a}$, that is
$|A|\le CN^{1-a/71}$: the baseline saving is $\eps=a/71$
(\cref{thm:record:lineardensity}). If the bound fails, let $\Lsc$ be the least
scale at which it fails, for some $A$. Then $(\Lsc,A)$ is a \emph{least
normalized counterexample} (\cref{def:record:counterexample}):
$Y=\Yf(\Lsc,A)>C_0\Lsc^{-a}\alpha$, while
$\Yf(\Lsc',A')\le C_0(\Lsc')^{-a}\alpha_{A'}$ for every scale
$\Lsc'<\Lsc$ and every square-difference-free $A'\subseteq[\Kint\Lsc']$,
$\alpha_{A'}$ its density. Every argument of the proof takes place at such a
pair, and its aim is a contradiction once $\Lsc$ is large. Throughout,
$o(1)$ is a quantity that tends to zero as $\Lsc\to\infty$, uniformly in $A$.

\paragraph{Fixing labels is rescaling.}
Let $\Blk$ be an atom of primes with $e_p=1$, and $q_\Blk=\prod_{p\in\Blk}p$.
Fixing, for one slot, its block label
$z_\Blk\in X_\Blk=\prod_{p\in\Blk}\Fp$ fixes for that slot a residue
class $\mathfrak a$ modulo $q_\Blk$ (the argument below is applied
slot by slot through the signed H\"older inequality): by the Chinese remainder theorem the specialized
lift $F_A(z_\Blk,\cdot)$ is, at every character $\xi$ of the other coordinates
with $q(\xi)q_\Blk\le N^{1/20}$, the lift of $q_\Blk\one_{A\cap(\mathfrak a+q_\Blk\Z)}$
(\cref{lem:record:fibre-scale}(a)). Split the class into the $q_\Blk$
progressions with the square step $q_\Blk^2$. Since $q_\Blk^2k^2=(q_\Blk k)^2$, the
set $A$ meets each progression $\{q_\Blk^2n'+c\}$ in a rescaled copy
$A'_c=\{n':q_\Blk^2n'+c\in A\}$, which is again square-difference-free and lives
at the scale $N/q_\Blk^2$. So, up to the characters of large conductor, whose
contribution is an error linear in $\alpha$ (\cref{lem:record:density-dc}),
$F_A(z_\Blk,\cdot)$ is an average of lifts of square-difference-free sets at a
shorter scale, composed with square-affine maps, which preserve $\mu$
(\cref{lem:record:fibre-scale}(b),(c)). By the triangle inequality for the
seminorm $\Lambda_\mu(\cdot)^{1/72}$ and the induction hypothesis at the shorter
scale, the specialized input has $\Lambda_\mu$ at most
$(1+o(1))(q_\Blk^2)^a(\alpha_{\rm f}/\alpha)Y$, $\alpha_{\rm f}$ the density of $A$ in the
class; the average over the progressions introduces no further
factor: by the triangle inequality for the seminorm
$\Lambda_\mu(\cdot)^{1/72}$ its seminorm is at most the weighted
average of the progression seminorms (the weights sum to one), and
concavity of $t\mapsto t^{1/72}$ moves that average inside the
power.
A change of the law at $\Blk$ changes all $72$ slots, and the signed H\"older
inequality multiplies the $72$ slot factors $(q_\Blk^2)^{a/72}$, while the
fibre densities average to $\alpha$ under labels with uniform one-slot
marginals.

\paragraph{The price of a change: total variation times $q_{\Blk}^{2a}$.}
The \emph{price} of a finite signed measure $\zeta$ on $X_\Blk^V$ is
$\|\zeta\|_{\rm var}\,q_\Blk^{2a}$ (\cref{def:price}; at a prefix prime the
factor is $p^{e_pa}$). The structured-energy estimate
(\cref{prop:record:density-structured}) states that, at a least normalized
counterexample, the \emph{energy} $\Lambda_{\zeta\otimes\mu_{-\Blk}}(F_A)$ of $\zeta$, in
which the factor $\mu_\Blk$ of the reference is replaced by $\zeta$, has absolute
value at most $(1+o(1))\|\zeta\|_{\rm var}\,q_\Blk^{2a}\,Y$, provided the variation of $\zeta$ has
uniform one-slot marginals; changes at several atoms multiply their prices.
Inputs restricted to intervals and residue classes are priced in the same
way. An input that sees only an interval of length $\Lsc=N/\Kint$ and a
residue class modulo an integer $D_0$ (the inputs $g^{I,D_0,s}$ of
\cref{subsec:mech:index}) is a lift at a scale shorter by the factor
$\Kint\sq(D_0)$, where $\sq(D_0)$ is the least square divisible by $D_0$ (for
instance $\sq(24)=144$). Its slot costs $(\Kint\sq(D_0))^{a/72}$, and a fixed
interval adds the factor $\Kint^{1/72}$, because the density of $A$ on an
interval of length $N/\Kint$ is at most $\Kint\alpha$
(\cref{def:record:structured-energy}). The product of these factors is the
\emph{coarse price} of the energy.

\paragraph{The only use of the induction hypothesis.}
Every comparison between laws in the proof is reduced to energies of this
kind, and the induction hypothesis enters only through the rescaling
above. Everything else is an exact identity at one scale, an estimate of a
finite local law, or one of two analytic remainders: the cancellation of
\cref{subsec:mech:kernel}, and the selections of atoms with a large product
(\cref{lem:record:density-can,lem:record:density-lp}). Both remainders are
$O(\Lsc^{-\sigma}\alpha)$, linear in the density. At a counterexample
$Y>C_0\Lsc^{-a}\alpha$, so such a remainder is at most
$(C/C_0)\Lsc^{a-\sigma}Y=o(1)Y$. This is why $a<\sigma$ is required, and why
the induction can keep the full density factor $\alpha$ on the right: the
density exponent $\theta=1$ of \cref{thm:record:lineardensity}.

\subsection{One query and its external index \texorpdfstring{$(b,s)$}{(b,s)}}\label{subsec:mech:index}

The prime $2$ and the primes $p\equiv3\pmod4$ below $Q_3$ form the fixed
external modulus
$\Mzero=2^{129}\prod_{3\le p<Q_3,\,p\equiv3\,(4)}p^{128}$.
They are not coordinates of $G_N$; their conductor depths are revealed
by an external random index attached to each query.

\paragraph{Intervals, inputs and the index.}
Split $[N]$ into the $\Kint$ consecutive intervals $I_b=((b-1)\Lsc,b\Lsc]\cap\Z$,
$b\in\Z/\Kint\Z$ (identified with $\{1,\dots,\Kint\}$). For an interval $I$, a
divisor $D_0$ of $\Mzero$ and $s\in\Z/D_0\Z$, the \emph{input} $g^{I,D_0,s}$ is
the lift of $D_0\one_{n\equiv s\,(D_0)}\one_A$ restricted to $I$, with
coefficients normalized by $|I|$ (\cref{def:record:functional}(f)); its mean
is the relative density $D_0\#\{n\in A\cap I:n\equiv s\ (D_0)\}/|I|$ of $A$ in
the class. Fix one query $e=(u_e,v_e)$. Its \emph{index} is $\omega=(b,s)$,
uniform on $\Z/\Kint\Z\times\Z/\Mzero\Z$, and its \emph{coupled index} is
$\omega'=(b+1,s+r)$, with $r$ independent of $\omega$ and distributed
as the square of a uniform residue conditioned on root valuation less than
$64$ in each prime-power coordinate. Thus $r$ is a nonvanishing square
class (\cref{def:nonvanishing-square-class,lem:imp:padic-envelope}). The
\emph{half-functions} $X_\omega$ and $Y_\omega$ carry the input $g^{I_b,\Mzero,s}$ at
the source $u_e$, respectively at the reflected target $\iota v_e$, and $F_A$
at the other slots of $V_+$ (\cref{def:hyb:coupled}; the half-function
$Y_\omega$ is unrelated to the average $Y$).

\paragraph{The mean identity.}
The lift is linear, the intervals $I_b$ partition $[N]$, and the residue
classes partition $\Z$; so averaging an input over a uniform index returns
the lift of $A$ (\cref{lem:twointerval}):
\begin{equation}\label{eq:mech:mean}
 \E_{b,s}\,g^{I_b,\Mzero,s}=F_A .
\end{equation}
Hence $\E_\omega X_\omega=\E_\omega Y_\omega=\bar X$, the half-function carrying $F_A$ at
every slot of $V_+$, and $Q_\kappa(\bar X,\bar X)=\Lambda_\kappa(F_A)=:Y_\kappa$ for every
measure $\kappa$; in particular $Y_\mu=Y$.

\paragraph{Coupling the two ends, and the signed combination $2P-T$.}
From now on cut forms are taken at the reserved cut $\iota$. Let $T$ be the
Markov operator $(Tf)(b,s)=\E_rf(b+1,s+r)$ on functions of the index
(this $T$ is an operator, not the dominator price $T$ of
Theorem~C), and $P$
the average over the whole index. For a finite signed measure $\kappa$ put
\[
 I_\kappa=\E_\omega Q_\kappa\bigl(X_\omega,(TY)_\omega\bigr),\qquad
 J_\kappa=\E_\omega Q_\kappa\bigl(X_\omega,((2P-T)Y)_\omega\bigr),
\]
the operators acting on the index of $Y_\omega$. Since $\iota(\iota v_e)=v_e$,
the \emph{coupled form} $I_\kappa$ is the average of $\Lambda_\kappa$ with
$g^{I_b,\Mzero,s}$ at $u_e$, $g^{I_{b+1},\Mzero,s+r}$ at $v_e$ and $F_A$ at the
other slots: it places the two ends of the query on the residue classes $s$
and $s+r$ of the adjacent intervals $I_b$ and $I_{b+1}$. By
\eqref{eq:mech:mean}, $(PY)_\omega=\bar X$, so the \emph{signed index form} is
\begin{equation}\label{eq:mech:signed}
 J_\kappa=2Y_\kappa-I_\kappa ,
\end{equation}
the case of one query of \eqref{eq:record:expansion}.

\paragraph{Decomposition by exact conductor.}
The Chinese remainder theorem splits the index into its interval coordinate
and prime-power coordinates. In the $p$-coordinate, conditional expectation
onto residues modulo $p^j$ gives a nested filtration; its orthogonal
differences are the exact-conductor sectors. Write $\delta_{p,j}$ for the
decreasing envelope on the norm of the shift at conductor $p^j$ proved in
\cref{lem:imp:padic-envelope}. Independently at the different primes choose
$J_p$ with $\Pr(J_p\ge j)=\delta_{p,j}$, and always reveal the interval.
On a nonconstant product sector of depths $\mathbf j$, both $T$ and $2P-T$
have norm at most $\prod_p\delta_{p,j_p}$, with the interval factor $1$.
Both operators fix constants. The coefficient $2$ used in the comparison
comes instead from the expansion $2P-T$: its fully averaged term is $2P$,
while the remaining term is the coupled form that cancels. There
is no contradiction in the two numbers: $1$ is the eigenvalue of
both operators on constant index functions, while $2$ is the
coefficient of the fully averaged term $Y_\kappa$ in the expansion
\eqref{eq:mech:signed}.

Let $\Jc=(J_p)_p$ and $D_\Jc=\prod_{p\mid\Mzero}p^{J_p}$. Conditional
expectation gives $X^\Jc=\E[X\mid b,s\bmod D_\Jc]$ and,
for any bracket matrix $M$ at the reserved cut (introduced
in the next paragraph),
$V_M(X)=\E_\Jc\E_\omega Q_M(X^\Jc_\omega,X^\Jc_\omega)$, and
likewise for $Y$.
The conditioned input is $g^{I_b,D_\Jc,s\bmod D_\Jc}$
(\cref{lem:twointerval}). The one affected slot in $X^\Jc$ and the one in
$Y^\Jc$ each contribute the slot factor
$(\Kint\sq(D_\Jc))^{a/72}$, hence $(\Kint\sq(D_\Jc))^{a/36}$
jointly after averaging their densities. Consequently
\begin{equation}\label{eq:mech:RI}
 R_I=\E_\Jc(\Kint\sq(D_\Jc))^{a/36}
 =M_{\rm pad}(\Kint)<1.247139907407.
\end{equation}
The numerical bound in \eqref{eq:mech:RI} is at the baseline
values $a=\PushA$ and $\Kint=\PushK$. The finite product formula
and its exact relation to the single-level square-shift moment are
proved in \cref{lem:hyb:moment,prop:imp:padic-price}.
Only the conditioning used to bound an energy has the random modulus
$D_\Jc$; the kernel cancellation itself always uses inputs at the
full modulus $\Mzero$.

\paragraph{The index-support lemma.}
Let $\bigl(\begin{smallmatrix}M_{\rm L}&\Gamma\\ \Gamma^{\mathsf T}&M_{\rm R}
\end{smallmatrix}\bigr)\succeq0$ be a \emph{bracket} of real matrices indexed by
$G_N^{V_+}$, and let $\mathcal T$ be an operator on functions of the index that
maps every exact-conductor sector into itself, is the identity on constants
and has norm at most $\prod_p\delta_{p,j_p}$ on the sector of depths
$\mathbf j$. Then \cref{lem:record:common-mean} gives
\begin{equation}\label{eq:mech:index-support}
 \bigl|\E_\omega\ip{X_\omega}{\Gamma(\mathcal TY)_\omega}\bigr|
 \le V_{\rm L}^{1/2}V_{\rm R}^{1/2},
\end{equation}
where $V_{\rm L}$ and $V_{\rm R}$ are the depth-conditioned energies under
$M_{\rm L}$ and $M_{\rm R}$. The probability that every $J_p\ge j_p$ equals
$\prod_p\delta_{p,j_p}$, exactly the weight of that sector in these energies.
The same lemma keeps the constant component exactly. With
$M_{\rm L}=M_{\rm R}=\Gamma\succeq0$ and $\mathcal T=T$ this is the plain
coupling inequality in the finite-depth form.

\subsection{Where square-difference-freeness enters}\label{subsec:mech:kernel}

In this one-query section a nonnegative law $\nu$ \emph{accepts}
the query $e$ if, $\nu$-almost surely, $z_{v_e}-z_{u_e}$ is a
nonzero square at every prime of $G_N$ with $e_p=1$, and (this is
the almost-sure condition alone; in Part~III \emph{accepts} also
includes the marginal requirements of (CP4), as in the glossary)
at every prime of the prefix it is nonzero, of even valuation $2j<e_p$, with
a square unit part (\cref{def:record:comparison-law}(a)). Consider the part
of the coupled form $I_\nu$ with $b\ne\Kint$. The two ends of the query carry
$g^{I_b,\Mzero,s}$ and $g^{I_{b+1},\Mzero,s+r}$. Every integer in the support
of the first, $A\cap I_b\cap(s+\Mzero\Z)$, is smaller than every integer in the
support of the second; their residues differ by $r$, a nonvanishing square class
modulo $\Mzero$; and at every prime of $G_N$ the law forces the labels of the
two ends to differ by a nonzero square. Fixing the labels at the primes
below $\Pfive=2^{65536}$ places the two ends in residue classes whose
difference is a square class of the kind the kernel admits
(\cref{lem:record:fibre-scale}(a)); at the primes $p\ge\Pfive$ the endpoint
pair is uniform on the strict-square pairs $\Om$, by the
\ref{it:record:R4} clause included in
\cref{def:record:comparison-law}(a) and not by the square condition
alone, and integrating it turns the average
into a Fourier pairing with the multiplier
$g(\xi)=\varphi(q(\xi))^{-1}\sum_{m\in(\Z/q(\xi)\Z)^\times}e(\xi m^2)$, the average
of $e(\xi m^2)$ over the units modulo $q(\xi)$ (\cref{lem:square-shift}(d)).

Three of the hypotheses of the actual-square kernel
(\cref{prop:record:kernel20}(iii)) are the form in which square-difference-freeness
is used: the two functions $F_1$ and $F_2$ are supported in single residue
classes $a_1$ and $a_2$ modulo $D$, with $a_2-a_1$ a nonvanishing square class;
every $x\in\supp F_1$ is smaller than every $y\in\supp F_2$; and no difference
$y-x$ is the square of a positive integer. Here $F_1$ and $F_2$ are built
from the square-difference-free set $A$, which is what makes the third
hypothesis true. Under these hypotheses the kernel bounds the Fourier
pairing $D\sum_{\xi\in\Xi_{H,D}}g(\xi)\widehat F_1(-\xi)\widehat F_2(\xi)$, truncated to
the frequencies $\Xi_{H,D}$ of squarefree conductor at most $H=N^{1/40}$ prime
to $D$, by $CH^{-1/3}\|F_1\|_{\ell^2}\|F_2\|_{\ell^2}/N$, where
$H^{-1/3}=N^{-1/120}$. In the proof
(\cref{app:kernel20}) this pairing is rewritten using a signed kernel whose
support consists of positive integer squares; the third hypothesis makes
that kernel's contribution exactly zero, and the factor $H^{-1/3}$ bounds
the remainder of the rewriting. With conductor decay at the large primes and
moment bounds for the other $70$ slots, the part of $I_\nu$ with $b\ne\Kint$ is
$O(\Lsc^{-\sigma}\alpha)=o(1)Y$ (\cref{lem:record:density-can}). This is the
only place where the absence of square differences yields a gain; the
rescaling of \cref{subsec:mech:prices} uses only that the property is
preserved.

If $b=\Kint$, the interval $I_\Kint$ is followed, cyclically, by $I_1$, and the
two ends are not ordered. This \emph{wrap} has probability $1/\Kint$; its term
is bounded by energies, and it becomes the term $W$ of the closing.

\subsection{One query: \texorpdfstring{$1\le\mathcal E+BR_I/2+W_1+o(1)$}{1 <= E+B R\_I/2+W\_1+o(1)} at a least counterexample}\label{subsec:mech:one-query}

This proposition is the whole mechanism with one query in place of
six: if the reference law is reflection positive, a comparison law
accepts the query and is dominated by a $2\times2$ block with energy
factor $B$, then a least counterexample would force
$1\le\mathcal E+BR_I/2+W_1+o(1)$. Its proof is
\cref{rem:cp:one-query}, the case $k=1$ of Theorem~C.

\begin{proposition}[One query]\label{prop:mech:one-query}
Consider the frame data of (CP0) (\cref{def:cp0}) with the index law of
\cref{subsec:mech:index} and a single query $e\in\mathcal E_6$, so $k=1$. Let
$0<a<\sigma$, let $(\Lsc,A)$ be a least normalized counterexample for
$(a,1,C_0)$ (\cref{def:record:counterexample}), and $Y=\Yf(\Lsc,A)$. Let the
reference $\mu$ satisfy (CP1) (\cref{def:cp1}): a product over the atoms of
probability measures, each reflection positive at the seven cuts of
$\mathfrak F$, square-affine invariant and mixing at the primes
$p\ge\Pfive$. Let $\nu=\bigotimes_\Blk\nu_\Blk$ be a product of nonnegative
measures, of total mass bounded independently of the scale, that accepts
$e$ in the sense of (CP4)(a) (\cref{def:record:comparison-law}; in particular
every one-slot marginal of $\nu_\Blk$ is a multiple of the uniform law).
Suppose that
\begin{enumerate}[label=\textup{(H\arabic*)}]
\item $Y_\nu\ge(1-\mathcal E-o(1))Y$ for a number $\mathcal E$; and
\item there are nonnegative measures $M_\Blk$ on $X_\Blk^V$ with
 $\bigl(\begin{smallmatrix}M_\Blk&\nu_\Blk\\ \nu_\Blk^*&M_\Blk\end{smallmatrix}\bigr)
 \succeq_\iota0$ for every atom, and a number $B$, such that the energy
 $\E_\omega\Lambda_M((g^{I_v,D_{0,v},s_v})_v)$ of every structured slot datum
 (\cref{def:record:structured-energy}(c); the expectation is
 over the index vector $\omega=(I_v,s_v)_v$ under the law
 $\mathbf P$ of that datum) under $M=\bigotimes_\Blk M_\Blk$ is at most
 $(B+o(1))$ times its coarse price times $Y$.
\end{enumerate}
Then, with the conductor moment $R_I$ of \eqref{eq:mech:RI},
\[
 1\le\mathcal E+\frac{BR_I}2+W_1+o(1),\qquad W_1=\frac{BR_I}{2\Kint^{35/36}} .
\]
\end{proposition}

\emph{Sketch of the proof.} By \eqref{eq:mech:signed}, $2Y_\nu=J_\nu+I_\nu$. The
brackets of (H2) tensorize over the atoms (\cref{lem:record:block-gram}), and
\eqref{eq:mech:index-support}, with $\Gamma$ the cut matrix of $\nu$ at $\iota$ and
$\mathcal T=2P-T$, gives $|J_\nu|\le V_{\rm L}^{1/2}V_{\rm R}^{1/2}$. By (H2) and the
slot prices of \cref{subsec:mech:index}, each of the depth-conditioned energies
$V_{\rm L}$ and $V_{\rm R}$ (the energies after conditioning on
the random conductor depths $\Jc$) is at most $(BR_I+o(1))Y$, the conductors being held fixed and averaged over $\Jc$
only at the end. In $I_\nu$ the part with $b\ne\Kint$ is $o(1)Y$ by the kernel
(\cref{subsec:mech:kernel}). The wrap $b=\Kint$ has probability $1/\Kint$, and
\eqref{eq:mech:index-support} on the residue coordinates bounds it by energies
at the fixed intervals $I_\Kint$ and $I_1$, which cost the density factor
$\Kint^{1/72}$ at each of the two changed slots: its term is at most
$(BR_I\Kint^{1/36}/\Kint+o(1))Y=(2W_1+o(1))Y$. With (H1),
$2(1-\mathcal E-o(1))Y\le2Y_\nu\le(BR_I+2W_1+o(1))Y$; divide by $2Y>0$. The
complete proof is \cref{rem:cp:one-query}, which checks the hypotheses of the
lemmas used in the frame of Part~III.

\emph{Keeping the common mean.} Part~(a) of \cref{lem:record:common-mean}
keeps exactly the component supported on no query and no conductor
coordinate (the \emph{empty-support}, that is constant, component):
that of $J_\nu$ is $Y_\nu$, and
that of each energy is the mean energy $e_M=Q_M(\bar X,\bar X)$ of the bracket
measure. It gives $J_\nu\le Y_\nu+(BR_I+o(1))Y-e_M$, and if
$e_M\ge(m-o(1))Y$, the conclusion becomes
\[
 1\le\tfrac12\mathcal E+\frac{BR_I}2+W_1+\frac{1-m}2+o(1):
\]
keeping the common mean halves the coefficient of $\mathcal E$, and a
mean constant $m>1$ lowers the right-hand side further through
$(1-m)/2$. This is the case $k=1$ of
\cref{thm:comparison} without a first stage ($L_0=0$), with the one energy
factor $B$ of (H2) in place of both $Be^T$ in $J$ and $Be^{S+A}$ in $W$.

\begin{remark}[A positive comparison law]\label{rem:mech:positive}
If every factor $\nu_\Blk$ were reflection positive at $\iota$, the bracket of
(H2) would hold with $M_\Blk=\nu_\Blk$, since
$\bigl(\begin{smallmatrix}C&C\\ C&C\end{smallmatrix}\bigr)=
\bigl(\begin{smallmatrix}1&1\\ 1&1\end{smallmatrix}\bigr)\otimes C\succeq0$ for
$C=C_\iota(\nu_\Blk)\succeq0$, and part (a) of the index-support lemma with
$\mathcal T=T$ would be the plain coupling inequality
$2Y_\nu-\bar V_\nu\le I_\nu$ of \cref{lem:hyb:coupling}, where
$\bar V_\nu=\frac12(V_\nu(X)+V_\nu(Y))$ is a scalar (not the slot
set $V$). The comparison laws of the
construction accept the queries but are not reflection positive at $\iota$
(\cref{subsec:walk:bracket} explains why); the signed index form and the
bracket are what replace positivity there. In the middle range the
comparison laws are signed measures: a fold record of the comparison law
is a signed combination of three nonnegative accepting laws with exact
marginals (\cref{def:record:mixture}). The cancellation of
\cref{subsec:mech:kernel} is applied to each nonnegative product component,
and the variation of the signed laws multiplies only the constant of a
vanishing remainder (\cref{def:cp4,lem:record:wrap-expansion}). The bracket asks less than
positivity: $M$ has to dominate $\nu$ only at the one cut $\iota$, and it need
not accept anything.
\end{remark}

\subsection{Why the sum of error bounds must be \texorpdfstring{$<1$}{<1}, and why one query fails}
\label{subsec:mech:budget}

\Cref{prop:mech:one-query} says that a least counterexample forces
$\mathcal E+BR_I/2+W_1\ge1-o(1)$. Hence if $\mathcal E$ and $B$ do
not depend on the scale and $\mathcal E+BR_I/2+W_1<1$, no least
normalized counterexample exists at large scales; once every $o(1)$
is shown to come only from structured energies and the two analytic
remainders (\cref{lem:record:rec}), the induction of
\cref{thm:record:lineardensity} gives $s(N)\le CN^{1-a/71}$. Each
term has a direct meaning. The left side is the constant coefficient $2$ of $2P-T$, divided by
$2$: in \eqref{eq:mech:signed} the term without an active index is $2Y_\nu$,
while the coupled form, in which the index is active, cancels by the kernel
except for the wrap. The conductor moment $R_I\ge1$ is the price of
revealing the index: the squared slot price of an input restricted to an
interval and to the residues kept by the inclusion law, averaged over the
inclusions. And $B$ is the energy factor of the bracket, the price of
controlling by a positive measure a law that is not reflection positive at
$\iota$; in the comparison theorem $B=e^{U_0}+R_{\rm err}\ge1$.

With one query the coupled term $BR_I/2$ is at least $R_I/2$, about
$0.623$ at the baseline $R_I$, so even with $B=1$ the comparison
price $\mathcal E$ would have to be below about $0.377$; at the
baseline energy factor $B<2.332628665480$ it is about $1.454>1$, so
one query cannot give $\Theta<1$. Using $k$ queries with disjoint
endpoints and independent indices, the operator
$\bigotimes_e(2P_e-T_e)$ has constant coefficient $2^k$, while the
energies factor over the queries, each contributing $R_I$
(\cref{lem:envelope-energy}). The coupled term becomes
\[
 J=Be^T\Bigl(\frac{R_I}2\Bigr)^k,
\]
where $T$ is the total price of the dominators for $\nu$ in (CP6)
(\cref{def:cp6}). For $k=6$, $(R_I/2)^6<0.0588$ at the baseline
$R_I$, and $J<0.25177$ at the baseline $B$ and $T<0.60760$
(baseline values; the improved values are in
\cref{sec:four-routes-closing}). Why six queries are used, and
why these six, is explained in \cref{subsec:walk:queries}.

Theorem~C is \cref{prop:mech:one-query} for $k$ queries with two
refinements: a first-stage law $\lambda$ with total price $L_0$
between $\mu$ and $\nu$, dominated at $\iota$ with energy factor $B$,
against which the symmetric and skew parts of $\nu-\lambda$ are
priced separately; and retention of the common mean, which lowers
the coefficient of $L_0+\mathcal E$ to $1-2^{-k}$. Its comparison
law may be, at a set of atoms whose number is bounded
independently of the scale (the middle blocks), a real combination
of
nonnegative accepting laws (\cref{def:cp4},
\cref{rem:mech:positive}). It gives $1\le\Theta+o(1)$ with
$\Theta=(1-2^{-k})(L_0+\mathcal E)+J+W+(1-m)/2^k$ at a least
counterexample. Part~V constructs $\nu$ with baseline
$\Theta<\PushTheta$ (\cref{sec:closing,tab:intro:values}); the
improved value $\Theta<\FourRoutesTheta$ is proved later by the
separate $\vartheta>1$ induction.

\subsection{Why conditioning independent labels on a square difference does not work}\label{subsec:mech:crude}

Which law should accept the query? The obvious candidate starts from the
law $\mu_0$ of independent uniform labels, whose cut matrices have rank one
and which is reflection positive at every cut, and conditions it on
$z_{v_e}-z_{u_e}\in\Qsq$ at every prime $p$ of $G_N$ with $e_p=1$ (and on the
corresponding event at the prefix primes). At one prime this replaces the
uniform law of the pair $(z_{u_e},z_{v_e})$ by the uniform law on the
strict-square pairs $\Om$. The uniform law puts the mass $(p+1)/(2p)$
outside $\Om$, and conditioning moves all of it onto $\Om$; so the change
$\nu_p-\mu_{0,p}$ has total variation norm
$\|\nu_p-\mu_{0,p}\|_{\rm var}=(p+1)/p>1$ (twice the moved mass),
and its price exceeds
$p^{2a}>1$. In the expansion of
$\nu=\bigotimes_p(\mu_{0,p}+(\nu_p-\mu_{0,p}))$ the prices of the primes
multiply, and their total grows with the number of primes of $G_N$: it is
not bounded independently of the scale, so $\mathcal E$ cannot be
bounded. Even a change of size about $1/p$ at every prime is too
large, since $\sum_pp^{2a-1}$ over the primes of $G_N$ is unbounded.

Two ingredients are missing. First, reference laws whose records
already give a square difference on most of their mass, so that an
accepting law changes only a small part at each prime. Second, a way
to bound many primes jointly: a change on a whole block of primes
should be bounded by about its worst prime, with one factor
$q_{\Blk}^{2a}$ for the block. \Cref{sec:design-walk} explains both
(\cref{subsec:walk:graph-laws,subsec:walk:completions} and
\cref{subsec:walk:blocks} respectively).

\paragraph{Consequences and what to read next.}
This section quotes, by their labels, statements of Part~II (the square
shift \cref{lem:square-shift}, \cref{sec:positivity}, the functional and the
kernel of \cref{def:record:functional,prop:record:kernel20}, and
\cref{sec:scale-comparison}) and of Part~III
(\cref{sec:comparison-theorem,sec:comparison-lemmas,sec:comparison-proof}), the
comparison law of the middle range (\cref{def:record:mixture}), and
the final values of \cref{sec:closing}; its only new statement,
\cref{prop:mech:one-query}, is proved in \cref{rem:cp:one-query}.
